\BeleskeID{09_2-s049}
% element: 09_2-s049:e05
Osnova rada CMOS kola je tanak oksid na gejtu. Ulaz preko gejta ima veoma veliku impedansu. Ako prima spoljašnji signal, bez zaštite nema putanju ni prema napajanju ni prema masi. Relativno mala količina naelektrisanja tada može podići napon i izazvati velika električna polja, što može dovesti do proboja oksida.

Zaštitne diode na ulaznim priključcima obezbeđuju odvođenje viška naelektrisanja prema napajanju ili prema masi. Njihova uloga u prikazanom objašnjenju je da spreče potencijale gejta iznad napajanja ili ispod mase. Zbog velike osetljivosti na statičko naelektrisanje, pinovi CMOS kola ne dodiruju se rukama. Pri rukovanju se oprema i operater povezuju sa potencijalom mase na kojem se nalazi i kolo.
